\documentclass[a4paper,11pt]{article}
\usepackage[utf8]{inputenc}
\usepackage[T1]{fontenc}
\usepackage[ngerman]{babel}
\usepackage{amsmath,amssymb}
\usepackage{xcolor}
\usepackage{colortbl}
\usepackage{array}
\usepackage{tikz}
\usepackage{enumitem}
\usepackage[margin=2cm]{geometry}
\definecolor{bgdark}{HTML}{1E1E1E}
\definecolor{fgtext}{HTML}{E6E6E6}
\definecolor{gridcol}{HTML}{4A4A4A}
\definecolor{accent}{HTML}{6CB6FF}
\definecolor{loescol}{HTML}{7FD18C}
\pagecolor{bgdark}
\color{fgtext}
\arrayrulecolor{fgtext}
\setlength{\parindent}{0pt}
\newcommand{\karo}[1]{\par\noindent\begin{tikzpicture}\draw[step=5mm,gridcol,very thin](0,0) grid (\linewidth,#1);\end{tikzpicture}\par\vspace{2mm}}
\newcommand{\titel}[1]{\begin{center}{\LARGE\color{accent}\textbf{Optimierungsverfahren}}\\[3pt]{\large #1}\end{center}\vspace{1mm}\noindent\textbf{Name:}\ \underline{\hspace{6cm}}\hfill\textbf{Datum:}\ \underline{\hspace{4cm}}\\[3pt]{\color{gridcol}\rule{\linewidth}{0.4pt}}\vspace{2mm}}
\newcommand{\thema}[2]{\clearpage{\large\color{accent}\textbf{Thema #1:\ #2}}\par\vspace{1mm}{\color{gridcol}\rule{\linewidth}{0.4pt}}\par\vspace{2mm}}
\newcommand{\bsp}{\par\vspace{1mm}\noindent{\color{loescol}\large\textbf{Musterbeispiel --- so rechnet man es}}\par\vspace{0.5mm}}
\newcommand{\loesweg}{\par\vspace{0.5mm}\noindent{\color{loescol}\textbf{Lösungsweg:}}\ }
\newcommand{\uebung}[2]{\clearpage\noindent{\large\color{accent}\textbf{Aufgabe #1}}\hfill{\color{gridcol}\normalsize #2}\par\vspace{1mm}}

\begin{document}
\titel{Großer Lern-Übungszettel --- Blatt 06: \textbf{Der komplette Stoff}}

\noindent Dieser Zettel deckt \textbf{alle} prüfungsrelevanten Themen der Optimierung ab. Zu \emph{jedem}
Thema gibt es:
\begin{itemize}[nosep,leftmargin=1.4em]
  \item ein \textbf{Musterbeispiel mit ausführlicher Erklärung} (so wird gerechnet), und
  \item \textbf{zwei Übungsaufgaben} zum selbst Rechnen (Karo-Platz darunter, Lösungen im Lösungs-PDF).
\end{itemize}
\vspace{1mm}
{\color{gridcol}\rule{\linewidth}{0.4pt}}

% ============================================================
% THEMA 1 — Modellierung
% ============================================================
\thema{1}{Modellierung eines Optimierungsproblems}
\bsp
Eine Firma stellt zwei Produkte $A$ und $B$ her. Pro Einheit beträgt der Gewinn $3$ (für $A$) bzw.\ $5$
(für $B$). Eine Einheit $A$ benötigt $1$\,h Maschinenzeit, $B$ benötigt $2$\,h; insgesamt stehen $40$\,h
zur Verfügung. Außerdem braucht $A$ $3$\,kg, $B$ $1$\,kg Material; vorhanden sind $60$\,kg.
Stelle das Optimierungsproblem auf.
\loesweg
\textbf{(1) Variablen:} $x_1=$ Stückzahl von $A$, $x_2=$ Stückzahl von $B$, beide $\ge 0$.
\textbf{(2) Zielfunktion} (Gewinn maximieren):
\[
\max\; z = 3x_1 + 5x_2 .
\]
\textbf{(3) Restriktionen} (jede Ressource $\le$ Vorrat):
\[
\underbrace{x_1 + 2x_2 \le 40}_{\text{Maschinenzeit}},\qquad
\underbrace{3x_1 + x_2 \le 60}_{\text{Material}},\qquad x_1,x_2\ge 0 .
\]
\emph{Merke:} immer (a) Entscheidungsvariablen benennen, (b) Ziel als $\min$/$\max$, (c) jede Ressource/Bedingung
als Ungleichung, (d) Nichtnegativität nicht vergessen.

\uebung{1.1 --- Mischungsproblem}{leicht}
Ein Landwirt mischt zwei Futtersorten $F_1$ ($0{,}30$\,€/kg) und $F_2$ ($0{,}50$\,€/kg). Pro kg liefert
$F_1$ $1$ Einheit Protein und $3$ Einheiten Vitamin, $F_2$ liefert $2$ Einheiten Protein und $1$ Einheit
Vitamin. Die Mischung soll \textbf{mindestens} $8$ Einheiten Protein und $6$ Einheiten Vitamin enthalten.
Modelliere das Problem (Kosten minimieren).
\karo{8cm}

\uebung{1.2 --- Nichtlineares Modell}{mittel}
Aus Karton soll eine \textbf{oben offene} Schachtel mit Kantenlängen $x,y,z>0$ gefertigt werden. Es steht
Material für eine Gesamt-Oberfläche von $12$ (Boden $+$ 4 Seitenwände) zur Verfügung. Modelliere das Problem
so, dass das \textbf{Volumen maximiert} wird. (Hinweis: Oberfläche $=xy+2xz+2yz$.)
\karo{8cm}

% ============================================================
% THEMA 2 — Konvexität
% ============================================================
\thema{2}{Konvexität von Funktionen und Mengen}
\bsp
Ist $f(x)=x^2-4x+3$ konvex? Ist die Menge $M=\{(x,y):x^2+y^2\le 4\}$ konvex?
\loesweg
\textbf{Funktion:} Eine $C^2$-Funktion einer Variablen ist konvex $\Leftrightarrow f''(x)\ge 0$ für alle $x$.
Hier $f'(x)=2x-4$, $f''(x)=2>0\ \Rightarrow$ \textbf{konvex} (sogar streng).
\par
\textbf{Mehrdim.:} $f$ konvex $\Leftrightarrow$ Hesse-Matrix $\mathbf H$ überall positiv semidefinit (alle Eigenwerte
$\ge 0$ bzw.\ alle führenden Hauptminoren $\ge 0$).
\par
\textbf{Menge:} $M$ konvex $\Leftrightarrow$ für je zwei Punkte $\vec a,\vec b\in M$ liegt die ganze Verbindungs\-strecke
in $M$. Eine Kreisscheibe (Norm-Kugel) erfüllt das $\Rightarrow$ $M$ ist \textbf{konvex}.

\uebung{2.1 --- Funktion zweier Variablen}{mittel}
Untersuche $f(x,y)=x^2+2y^2-xy$ auf Konvexität (Hesse-Matrix aufstellen, Definitheit über führende
Hauptminoren oder Eigenwerte bestimmen).
\karo{8cm}

\uebung{2.2 --- Zwei Kurzfragen}{leicht--mittel}
\begin{enumerate}[label=(\alph*),nosep]
  \item Ist $g(x,y)=x^2-y^2$ konvex, konkav oder keines von beidem? Begründe über die Hesse-Matrix.
  \item Ist $h(x)=x^3$ auf ganz $\mathbb R$ konvex? Begründe.
\end{enumerate}
\karo{8cm}

% ============================================================
% THEMA 3 — Analytische Optima (1D)
% ============================================================
\thema{3}{Analytische Optima --- eine Variable}
\bsp
Bestimme die lokalen Extrema von $f(x)=x^3-6x^2+9x+1$.
\loesweg
\textbf{(1) Notwendige Bedingung} $f'(x)=0$:
\[
f'(x)=3x^2-12x+9=3(x-1)(x-3)=0 \;\Rightarrow\; x=1,\ x=3 .
\]
\textbf{(2) Hinreichende Bedingung} über $f''(x)=6x-12$:
\[
f''(1)=-6<0 \Rightarrow \text{Maximum},\quad f(1)=1-6+9+1=5;\qquad
f''(3)=6>0 \Rightarrow \text{Minimum},\quad f(3)=1 .
\]
\emph{Ergebnis:} lokales Maximum $(1,5)$, lokales Minimum $(3,1)$.

\uebung{3.1 --- Polynom 4.\ Grades}{mittel}
Bestimme alle lokalen Extrema von $f(x)=x^4-8x^2+5$ und gib an, welches das globale Minimum ist.
\karo{8cm}

\uebung{3.2 --- Mit e-Funktion}{mittel}
Bestimme das Minimum von $f(x)=e^{x}-3x$. (Funktionswert auf 3 Nachkommastellen angeben.)
\karo{8cm}

% ============================================================
% THEMA 4 — Analytische Optima (2D)
% ============================================================
\thema{4}{Analytische Optima --- zwei Variablen (Hesse)}
\bsp
Bestimme das Optimum von $f(x,y)=x^2+xy+y^2-3x$.
\loesweg
\textbf{(1) Gradient $=\vec 0$:}
\[
\nabla f=\begin{pmatrix}2x+y-3\\ x+2y\end{pmatrix}\overset{!}{=}\vec 0 .
\]
Aus der 2.\ Gleichung $x=-2y$, eingesetzt: $2(-2y)+y-3=0\Rightarrow -3y=3\Rightarrow y=-1,\ x=2$.
\par\textbf{(2) Hesse-Matrix:}
$\mathbf H=\begin{pmatrix}2&1\\1&2\end{pmatrix}$, Hauptminoren $2>0$, $\det=3>0\Rightarrow$ positiv definit
$\Rightarrow$ \textbf{Minimum}. $f(2,-1)=4-2+1-6=-3$.

\uebung{4.1 --- Mit Sattelpunkt}{schwer}
Bestimme die stationären Punkte von $f(x,y)=x^3+y^3-3xy$ und klassifiziere sie (Hesse-Matrix).
\karo{8.5cm}

\uebung{4.2 --- Klassifikation}{mittel}
Bestimme den stationären Punkt von $f(x,y)=x^2-4xy+y^2$ und gib über die Hesse-Matrix an, ob es sich um
Minimum, Maximum oder Sattelpunkt handelt.
\karo{7cm}

% ============================================================
% THEMA 5 — Grafische LP + Slack
% ============================================================
\thema{5}{Lineare Optimierung --- grafisch \& Slackvariablen}
\bsp
$\max z=3x_1+2x_2$ unter $x_1+x_2\le 4$, $x_1+3x_2\le 6$, $x_1,x_2\ge 0$.
\loesweg
\textbf{(1) Zulässige Ecken} (Schnittpunkte der Restriktionsgeraden \& Achsen):
$(0,0),(4,0),(0,2)$ und Schnitt $x_1+x_2=4,\ x_1+3x_2=6$: Subtraktion $\Rightarrow 2x_2=2\Rightarrow x_2=1,x_1=3$,
also $(3,1)$.
\par\textbf{(2) Zielwert in jeder Ecke:} $(4,0)\!:12$, $(3,1)\!:11$, $(0,2)\!:4$ $\Rightarrow$ Optimum $(4,0)$, $z=12$.
\par\textbf{(3) Slackvariablen} $s_i=$ rechte Seite $-$ linke Seite im Optimum:
$s_1=4-(4+0)=0$ (Restriktion 1 aktiv), $s_2=6-(4+0)=2$ (Restriktion 2 hat $2$ Reserve).

\uebung{5.1 --- Grafisch lösen}{mittel}
$\max z=x_1+x_2$ unter $2x_1+x_2\le 10$, $x_1+3x_2\le 15$, $x_1,x_2\ge 0$. Bestimme das Optimum (Ecken
vergleichen) und die Werte der Slackvariablen im Optimum.
\karo{9cm}

\uebung{5.2 --- Bruch-Optimum}{mittel}
$\max z=4x_1+3x_2$ unter $2x_1+x_2\le 8$, $x_1+2x_2\le 8$, $x_1,x_2\ge 0$. Bestimme das Optimum und die
Slackvariablen.
\karo{9cm}

% ============================================================
% THEMA 6 — Standardform
% ============================================================
\thema{6}{Lineare Optimierung --- Standardform}
\bsp
Bringe in Standardform ($\min$, nur Gleichungen, alle Variablen $\ge 0$):
$\max z=2x_1+3x_2$ unter $x_1+x_2\le 4$, $x_1-x_2\ge 1$, $x_1\ge 0$, $x_2$ \emph{frei}.
\loesweg
\textbf{(1)} $\max z \Leftrightarrow \min(-z)=\min(-2x_1-3x_2)$.\quad
\textbf{(2)} $\le$ mit \emph{Slack} $s_1\ge0$, $\ge$ mit \emph{Surplus} $s_2\ge0$.\quad
\textbf{(3)} freie Variable splitten: $x_2=x_2^+-x_2^-$ mit $x_2^+,x_2^-\ge 0$.
\[
\min\; -2x_1-3x_2^++3x_2^-\quad\text{u.\,d.\,N.}
\]
\[
x_1+x_2^+-x_2^-+s_1=4,\qquad x_1-x_2^++x_2^--s_2=1,\qquad x_1,x_2^+,x_2^-,s_1,s_2\ge 0 .
\]

\uebung{6.1 --- Standardform I}{leicht}
Bringe in Standardform: $\max z=5x_1+4x_2$ unter $6x_1+4x_2\le 24$, $x_1+2x_2\le 6$, $x_1,x_2\ge 0$.
\karo{7cm}

\uebung{6.2 --- Standardform II}{mittel}
Bringe in Standardform: $\min z=3x_1+x_2$ unter $x_1+x_2\ge 3$, $2x_1+x_2\le 8$, $x_1$ \emph{frei}, $x_2\ge 0$.
\karo{8cm}

% ============================================================
% THEMA 7 — Ganzzahlige Optimierung
% ============================================================
\thema{7}{Ganzzahlige Optimierung (grafisch)}
\bsp
$\max z=x_1+x_2$ unter $3x_1+2x_2\le 12$, $x_1+2x_2\le 6$, $x_1,x_2\in\mathbb Z_{\ge 0}$.
\loesweg
\textbf{(1) Kontinuierliches Optimum:} Schnitt $3x_1+2x_2=12$, $x_1+2x_2=6\Rightarrow 2x_1=6\Rightarrow x_1=3,\ x_2=1{,}5$,
$z=4{,}5$.
\par\textbf{(2) Ganzzahlig:} runde/prüfe Nachbarn. $(3,1)$: $11\le12,\ 5\le6$ zulässig, $z=4$. $(2,2)$: $10\le12,\ 6\le6$
zulässig, $z=4$. Ein zulässiger ganzzahliger Punkt mit $z=5$ existiert nicht.
\par\textbf{Ergebnis:} $\vec x^*_{\mathbb Z}=(3,1)$ bzw.\ $(2,2)$, $z=4$ ($<4{,}5$: das kont.\ Optimum ist obere Schranke).

\uebung{7.1 --- Ganzzahlig I}{mittel}
$\max z=4x_1+5x_2$ unter $x_1+4x_2\le 10$, $3x_1+2x_2\le 12$, $x_1,x_2\in\mathbb Z_{\ge0}$. Bestimme das
kontinuierliche und das ganzzahlige Optimum.
\karo{9cm}

\uebung{7.2 --- Ganzzahlig II}{mittel}
$\max z=5x_1+4x_2$ unter $6x_1+4x_2\le 24$, $x_1+2x_2\le 6$, $x_1,x_2\in\mathbb Z_{\ge0}$. Bestimme das
kontinuierliche und das ganzzahlige Optimum.
\karo{9cm}

% ============================================================
% THEMA 8 — Transport: Nordwesteckenregel
% ============================================================
\thema{8}{Transportproblem --- Nordwesteckenregel}
\bsp
Angebot $a=(20,30,25)$, Nachfrage $b=(10,25,40)$, Kostenmatrix
$C=\begin{pmatrix}8&6&10\\9&12&7\\3&5&4\end{pmatrix}$. Bestimme eine Startlösung mit der Nordwesteckenregel und
ihre Kosten.
\loesweg
Starte oben links, belege jede Zelle maximal, gehe dann nach rechts/unten (Kosten werden \emph{ignoriert}):
\[
\begin{array}{l}
x_{11}=\min(20,10)=10\ (\text{Nachfr.\,1 fertig, Ang.\,1}\to10)\\
x_{12}=\min(10,25)=10\ (\text{Ang.\,1 fertig, Nachfr.\,2}\to15)\\
x_{22}=\min(30,15)=15\ (\text{Nachfr.\,2 fertig, Ang.\,2}\to15)\\
x_{23}=\min(15,40)=15\ (\text{Ang.\,2 fertig, Nachfr.\,3}\to25)\\
x_{33}=\min(25,25)=25
\end{array}
\quad
X=\begin{pmatrix}10&10&0\\0&15&15\\0&0&25\end{pmatrix}
\]
Basiszellen $=5=m+n-1$ \checkmark. Kosten $=10\cdot8+10\cdot6+15\cdot12+15\cdot7+25\cdot4=80+60+180+105+100=\boxed{525}$.

\uebung{8.1 --- NWE (2$\times$3)}{leicht}
Angebot $(30,30)$, Nachfrage $(20,25,15)$, $C=\begin{pmatrix}3&5&7\\6&4&2\end{pmatrix}$. Bestimme die
NWE-Startlösung und ihre Kosten.
\karo{7.5cm}

\uebung{8.2 --- NWE (2$\times$3)}{leicht}
Angebot $(40,30)$, Nachfrage $(25,30,15)$, $C=\begin{pmatrix}7&3&12\\2&5&6\end{pmatrix}$. Bestimme die
NWE-Startlösung und ihre Kosten.
\karo{7.5cm}

% ============================================================
% THEMA 9 — Transport: Matrix-Minimum
% ============================================================
\thema{9}{Transportproblem --- Matrix-Minimum-Methode}
\bsp
Gleiche Daten wie Thema 8: $a=(20,30,25)$, $b=(10,25,40)$, $C=\begin{pmatrix}8&6&10\\9&12&7\\3&5&4\end{pmatrix}$.
\loesweg
Wähle stets die noch freie Zelle mit den \emph{kleinsten} Kosten und belege sie maximal:
\[
\begin{array}{ll}
c_{31}=3: & x_{31}=\min(25,10)=10\ \Rightarrow \text{Spalte 1 fertig, Ang.\,3}\to15\\
c_{33}=4: & x_{33}=\min(15,40)=15\ \Rightarrow \text{Ang.\,3 fertig, Nachfr.\,3}\to25\\
c_{12}=6: & x_{12}=\min(20,25)=20\ \Rightarrow \text{Ang.\,1 fertig, Nachfr.\,2}\to5\\
c_{23}=7: & x_{23}=\min(30,25)=25\ \Rightarrow \text{Nachfr.\,3 fertig, Ang.\,2}\to5\\
c_{22}=12: & x_{22}=\min(5,5)=5
\end{array}
\quad
X=\begin{pmatrix}0&20&0\\0&5&25\\10&0&15\end{pmatrix}
\]
Kosten $=10\cdot3+15\cdot4+20\cdot6+25\cdot7+5\cdot12=30+60+120+175+60=\boxed{445}$ (besser als NWE $525$).

\uebung{9.1 --- MMM I}{mittel}
Angebot $(20,30)$, Nachfrage $(10,25,15)$, $C=\begin{pmatrix}4&6&8\\5&3&7\end{pmatrix}$. Bestimme die
Startlösung per Matrix-Minimum-Methode und die Kosten.
\karo{8cm}

\uebung{9.2 --- MMM II}{mittel}
Angebot $(50,40)$, Nachfrage $(30,50,10)$, $C=\begin{pmatrix}4&6&8\\5&3&7\end{pmatrix}$. Bestimme die
Startlösung per Matrix-Minimum-Methode und die Kosten.
\karo{8cm}

% ============================================================
% THEMA 10 — Transport: Vogel
% ============================================================
\thema{10}{Transportproblem --- Vogelsche Approximation (VAM)}
\bsp
Gleiche Daten wie Thema 8: $a=(20,30,25)$, $b=(10,25,40)$, $C=\begin{pmatrix}8&6&10\\9&12&7\\3&5&4\end{pmatrix}$.
\loesweg
Penalty $=$ Differenz der zwei kleinsten Kosten je Zeile/Spalte; größte Penalty wählen, dort günstigste Zelle
maximal belegen.
\[
\begin{array}{l}
\textbf{S1:}\ \text{Pen Z}=(2,2,1),\ \text{Pen S}=(5,1,3).\ \text{Max}=5\,(\text{Sp.\,1}),\ c_{31}=3:\ x_{31}=\min(25,10)=10.\\
\textbf{S2:}\ \text{Pen Z}=(4,5),\ \text{Pen S}=(1,3)\ [\text{Sp.\,2,3}].\ \text{Max}=5\,(\text{Z.\,2}),\ c_{23}=7:\ x_{23}=\min(30,40)=30.\\
\textbf{S3:}\ \text{Pen Z}=(4,-),\ \text{Pen S}=(1,6)\ [\text{Sp.\,2,3}].\ \text{Max}=6\,(\text{Sp.\,3}),\ c_{33}=4:\ x_{33}=\min(15,10)=10.\\
\textbf{S4:}\ \text{Rest Spalte 2 (Nachfr.\,25):}\ x_{12}=\min(20,25)=20,\ x_{32}=5.
\end{array}
\]
\[
X=\begin{pmatrix}0&20&0\\0&0&30\\10&5&10\end{pmatrix},\quad
\text{Kosten}=10\cdot3+30\cdot7+10\cdot4+20\cdot6+5\cdot5=30+210+40+120+25=\boxed{425}.
\]
Vergleich: NWE $525$, MMM $445$, \textbf{VAM $425$}.

\uebung{10.1 --- VAM I}{mittel}
Angebot $(30,30)$, Nachfrage $(20,25,15)$, $C=\begin{pmatrix}3&5&7\\6&4&2\end{pmatrix}$. Bestimme die
Startlösung mit der Vogelschen Approximationsmethode und die Kosten.
\karo{9cm}

\uebung{10.2 --- VAM II}{schwer}
Angebot $(40,30)$, Nachfrage $(25,30,15)$, $C=\begin{pmatrix}7&3&12\\2&5&6\end{pmatrix}$. Bestimme die
VAM-Startlösung und die Kosten. Vergleiche mit der NWE-Lösung aus Aufgabe 8.2.
\karo{9cm}

% ============================================================
% THEMA 11 — Transport: Optimalitätsprüfung (MODI)
% ============================================================
\thema{11}{Transportproblem --- Optimalität prüfen (MODI/Stepping-Stone)}
\bsp
Angebot $(50,40)$, Nachfrage $(30,50,10)$, $C=\begin{pmatrix}4&6&8\\5&3&7\end{pmatrix}$. Prüfe die
NWE-Startlösung auf Optimalität und verbessere sie ggf.
\loesweg
\textbf{NWE-Start:} $x_{11}=30,x_{12}=20,x_{22}=30,x_{23}=10$. Kosten $=30\cdot4+20\cdot6+30\cdot3+10\cdot7=400$.
\par\textbf{MODI:} setze $u_1=0$, dann $u_i+v_j=c_{ij}$ für Basiszellen:
$v_1=4,\ v_2=6,\ u_2=c_{22}-v_2=3-6=-3,\ v_3=c_{23}-u_2=7+3=10$.
\par Reduzierte Kosten $d_{ij}=c_{ij}-(u_i+v_j)$ der Nichtbasiszellen:
\[
d_{13}=8-(0+10)=-2<0,\qquad d_{21}=5-(-3+4)=4\ge0 .
\]
$d_{13}<0\Rightarrow$ \textbf{nicht optimal}. Zelle $(1,3)$ aufnehmen, Zyklus $(1,3)^+\!\to(2,3)^-\!\to(2,2)^+\!\to(1,2)^-$;
$\theta=\min(x_{23},x_{12})=\min(10,20)=10$.
\par Neu: $x_{13}=10,x_{23}=0,x_{22}=40,x_{12}=10$. Kosten $=400-10\cdot2=380$.
\par Erneuter MODI-Test: alle $d_{ij}\ge0$ $\Rightarrow$ \textbf{optimal}, Gesamtkosten $\boxed{380}$.

\uebung{11.1 --- Optimal?}{mittel}
Angebot $(30,30)$, Nachfrage $(20,25,15)$, $C=\begin{pmatrix}3&5&7\\6&4&2\end{pmatrix}$. Bestimme die
NWE-Startlösung und prüfe per MODI, ob sie bereits optimal ist.
\karo{9.5cm}

\uebung{11.2 --- Eine Verbesserung}{schwer}
Angebot $(40,30)$, Nachfrage $(25,30,15)$, $C=\begin{pmatrix}7&3&12\\2&5&6\end{pmatrix}$. Bestimme die
NWE-Startlösung, prüfe per MODI und führe (falls nötig) genau einen Verbesserungsschritt durch.
\karo{9.5cm}

% ============================================================
% THEMA 12 — Gradientenabstieg
% ============================================================
\thema{12}{Nichtlinear --- Gradientenabstieg}
\bsp
$f(x,y)=x^2+2y^2$, Startpunkt $\vec x_0=(2,1)$, Schrittweite $\lambda=0{,}1$. Führe einen Schritt aus.
\loesweg
$\nabla f=\begin{pmatrix}2x\\4y\end{pmatrix}\Rightarrow \nabla f(2,1)=\begin{pmatrix}4\\4\end{pmatrix}$.
\[
\vec x_1=\vec x_0-\lambda\nabla f(\vec x_0)=\begin{pmatrix}2\\1\end{pmatrix}-0{,}1\begin{pmatrix}4\\4\end{pmatrix}
=\begin{pmatrix}1{,}6\\0{,}6\end{pmatrix}.
\]
$f(\vec x_0)=6\ \to\ f(\vec x_1)=1{,}6^2+2\cdot0{,}6^2=2{,}56+0{,}72=3{,}28$ (Abnahme \checkmark).

\uebung{12.1 --- Ein Schritt 2D}{leicht}
$f(x,y)=x^2+y^2+xy$, Start $(1,1)$, $\lambda=0{,}1$. Berechne $\vec x_1$ und $f(\vec x_1)$.
\karo{7cm}

\uebung{12.2 --- Zwei Schritte 1D}{mittel}
$f(x)=(x-2)^2$, Start $x_0=0$, $\lambda=0{,}25$. Führe \textbf{zwei} Schritte aus und gib $x_1,x_2$ an.
\karo{7cm}

% ============================================================
% THEMA 13 — Newton-Verfahren
% ============================================================
\thema{13}{Nichtlinear --- Newton-Verfahren}
\bsp
Minimiere $f(x)=x^3-2x+2$ mit dem Newton-Verfahren ab $x_0=1$ (ein Schritt).
\loesweg
Newton zur Optimierung: $x_{k+1}=x_k-\dfrac{f'(x_k)}{f''(x_k)}$ mit $f'(x)=3x^2-2$, $f''(x)=6x$.
\[
x_1=1-\frac{f'(1)}{f''(1)}=1-\frac{3-2}{6}=1-\tfrac16\approx 0{,}8333 .
\]
(Konvergiert gegen $x^*=\sqrt{2/3}\approx0{,}8165$, das Minimum.) Mehrdim.: $\vec x_{k+1}=\vec x_k-\mathbf H^{-1}\nabla f$.

\uebung{13.1 --- Zwei Schritte 1D}{mittel}
$f(x)=x^3-6x^2+9x$. Führe ab $x_0=0$ \textbf{zwei} Newton-Schritte aus. Gegen welchen stationären Punkt
konvergiert das Verfahren?
\karo{8cm}

\uebung{13.2 --- Newton in 2D}{schwer}
$f(x,y)=x^2+xy+y^2$, Start $\vec x_0=(1,1)$. Führe einen Newton-Schritt
$\vec x_1=\vec x_0-\mathbf H^{-1}\nabla f(\vec x_0)$ aus. Was fällt auf?
\karo{8.5cm}

% ============================================================
% THEMA 14 — KKT
% ============================================================
\thema{14}{Nichtlinear --- KKT-Bedingungen}
\bsp
$\min f(x,y)=(x-2)^2+(y-1)^2$ unter $g(x,y)=x+y-1\le 0$.
\loesweg
\textbf{KKT:} $\nabla f+\mu\nabla g=\vec 0$, $\ \mu\ge0$, $\ \mu\,g=0$, $\ g\le0$.
$\nabla f=(2(x-2),2(y-1))$, $\nabla g=(1,1)$:
\[
2(x-2)+\mu=0,\qquad 2(y-1)+\mu=0\ \Rightarrow\ x-2=y-1\ \Rightarrow\ x=y+1 .
\]
\textbf{Fall $\mu=0$:} $x=2,y=1\Rightarrow g=2>0$ unzulässig.
\textbf{Fall $g=0$:} $x+y=1$ und $x=y+1\Rightarrow 2y=0\Rightarrow y=0,x=1$; $\mu=-2(x-2)=2\ge0$ \checkmark.
\par\emph{Ergebnis:} $\vec x^*=(1,0)$, $\mu=2$, $f=2$.

\uebung{14.1 --- KKT mit $\ge$-Restriktion}{mittel}
$\min f(x,y)=x^2+y^2$ unter $x+y\ge 4$. Stelle die KKT-Bedingungen auf und löse sie vollständig.
\karo{9cm}

\uebung{14.2 --- KKT mit Kreis}{schwer}
$\min f(x,y)=(x-1)^2+(y-2)^2$ unter $x^2+y^2\le 1$. Bestimme den KKT-Punkt, den Multiplikator $\mu$ und
$f$ im Optimum.
\karo{9.5cm}

% ============================================================
% THEMA 15 — Grafische nichtlineare Optimierung
% ============================================================
\thema{15}{Nichtlinear --- grafische Optimierung}
\bsp
$\min f(x,y)=(x-3)^2+(y-2)^2$ unter $x+y\le 4$, $x,y\ge0$.
\loesweg
$f$ misst den quadrierten Abstand zum Punkt $(3,2)$; die Höhenlinien sind \emph{Kreise} um $(3,2)$. Das
unbeschränkte Minimum $(3,2)$ verletzt $x+y\le4$ ($3+2=5$). Gesucht ist also der \emph{nächstgelegene
zulässige Punkt} = Lotfußpunkt von $(3,2)$ auf die Gerade $x+y=4$:
\[
\vec x^*=(3,2)-\frac{(3+2)-4}{2}(1,1)=(3,2)-\tfrac12(1,1)=(2{,}5,\,1{,}5),\quad f=0{,}25+0{,}25=0{,}5 .
\]
(Optimum dort, wo ein Höhenkreis die Restriktionsgerade gerade berührt.)

\uebung{15.1 --- Produkt maximieren}{mittel}
$\max f(x,y)=x\,y$ unter $x+y=10$, $x,y\ge0$. Bestimme das Optimum grafisch/analytisch (Höhenlinien sind
Hyperbeln).
\karo{7.5cm}

\uebung{15.2 --- Abstand zu einer Kreisscheibe}{schwer}
$\min f(x,y)=x^2+y^2$ unter $(x-2)^2+(y-2)^2\le 1$. Bestimme den optimalen Punkt und $f$ (Höhenlinien sind
Kreise um den Ursprung).
\karo{8.5cm}

% ============================================================
% THEMA 16 — Nelder-Mead
% ============================================================
\thema{16}{Gradientenfrei --- Nelder-Mead (Simplex)}
\bsp
$f(x,y)=x^2+y^2$. Simplex $B=(1,1),\ G=(2,2),\ W=(3,1)$. Führe eine Iteration aus.
\loesweg
\textbf{(1) Sortieren} nach $f$: $f(B)=2\le f(G)=8\le f(W)=10$ ($B$ best, $G$ zweitschlecht., $W$ schlecht.).
\par\textbf{(2) Schwerpunkt} $M$ aller Punkte \emph{außer} $W$: $M=\tfrac12(B+G)=(1{,}5,\,1{,}5)$.
\par\textbf{(3) Reflexion} $R=M+(M-W)=2M-W=(3,3)-(3,1)=(0,2)$, $f(R)=4$.
\par\textbf{(4) Entscheidung:} Regel ist
\begin{itemize}[nosep,leftmargin=1.4em]
  \item $f(R)<f(B)$: Expansion versuchen;
  \item $f(B)\le f(R)<f(G)$: $R$ \emph{annehmen};
  \item $f(R)\ge f(G)$: Kontraktion; bringt auch das nichts: Schrumpfung.
\end{itemize}
Hier $f(B)=2\le f(R)=4<f(G)=8\Rightarrow$ \textbf{$R$ annehmen}; $W$ wird durch $(0,2)$ ersetzt.
\par\textbf{Folgeschritt} (neuer Simplex $\{(1,1),(0,2),(2,2)\}$, neuer Worst $=(2,2)$): $M=\tfrac12((1,1)+(0,2))=(0{,}5,1{,}5)$,
$R'=2M-(2,2)=(-1,1)$, $f(R')=2$; wegen $f(B)\le f(R')<f(G)$ erneut annehmen.

\uebung{16.1 --- Expansion}{mittel}
$f(x,y)=x^2+2y^2$. Simplex $P_1=(2,0),\ P_2=(0,2),\ P_3=(2,2)$. Führe eine Iteration (Reflexion und ggf.\
Expansion) durch.
\karo{8.5cm}

\uebung{16.2 --- Kontraktion}{schwer}
$f(x,y)=x^2+y^2$. Simplex $P_1=(0,1),\ P_2=(1,0),\ P_3=(2,2)$. Führe eine Iteration durch (was passiert,
wenn die Reflexion schlechter als der zweitschlechteste Punkt ist?).
\karo{8.5cm}

% ============================================================
% THEMA 17 — DIRECT
% ============================================================
\thema{17}{Gradientenfrei --- DIRECT}
\bsp
Minimiere $f(x)=(x-0{,}3)^2$ auf $[0,1]$ mit DIRECT (eine Unterteilung).
\loesweg
DIRECT wertet stets die \emph{Mitte} eines Intervalls aus und drittelt das/die \emph{potentiell optimalen}
Intervalle. Start: Mitte $0{,}5$, $f=0{,}04$, Länge $1$. Dritteln in $[0,\tfrac13],[\tfrac13,\tfrac23],[\tfrac23,1]$
mit Mitten $\tfrac16,\tfrac12,\tfrac56$:
\[
f(\tfrac16)=(0{,}1\overline{6}-0{,}3)^2\approx0{,}0178,\quad f(\tfrac12)=0{,}04,\quad f(\tfrac56)\approx0{,}284 .
\]
Alle drei Intervalle sind gleich groß $\Rightarrow$ potentiell optimal ist das mit dem \emph{kleinsten} $f$,
also $[0,\tfrac13]$ (Mitte $\tfrac16$). Dieses wird als Nächstes gedrittelt.

\uebung{17.1 --- DIRECT 1D}{mittel}
Minimiere $f(x)=(x-0{,}7)^2$ auf $[0,1]$: führe die erste Unterteilung durch, berechne die drei Mittenwerte
und gib das als Nächstes zu unterteilende Intervall an.
\karo{7.5cm}

\uebung{17.2 --- Potentiell optimale Rechtecke}{schwer}
In der $(d,f)$-Darstellung (Größe $d$, Funktionswert $f$ der Mitte) liegen die Rechtecke
\[
\text{Nr.1}(1,6),\ \text{Nr.2}(1,4),\ \text{Nr.3}(3,5),\ \text{Nr.4}(3,3),\ \text{Nr.5}(9,4).
\]
Welche Rechtecke sind \emph{potentiell optimal}? Begründe (je Größe nur das beste, dann untere-rechte Hülle).
\karo{8cm}

% ============================================================
% THEMA 18 — (1+1)-ES
% ============================================================
\thema{18}{Metaheuristiken --- $(1+1)$-ES mit 1/5-Erfolgsregel}
\bsp
$f(x)=x^2$, Start $x=8$, $\sigma=4$. Regel: Erfolgsrate $s=\frac{\#\text{Erfolge}}{\text{Iterationen}}$;
$s>\tfrac15\Rightarrow\sigma\cdot2$, $s<\tfrac15\Rightarrow\sigma/2$, $s=\tfrac15\Rightarrow$ unverändert. Mutation
$x^{\text{neu}}=x+\sigma\cdot z$, akzeptiert nur bei Verbesserung. $z$-Folge: $1,\,-3,\,-0{,}5,\,0{,}5,\,-0{,}1$.
\loesweg
\[
\renewcommand{\arraystretch}{1.25}
\begin{array}{c|c|c|c|c|c}
\text{It} & x^{\text{neu}}=x+\sigma z & f^{\text{neu}} & \text{Erg.} & s & \sigma_{\text{neu}}\\\hline
1 & 8+4(1)=12 & 144>64 & \text{Miss} & 0/1=0 & 4/2=2\\
2 & 8+2(-3)=2 & 4<64 & \text{Erfolg} & 1/2=0{,}5 & 2\cdot2=4\\
3 & 2+4(-0{,}5)=0 & 0<4 & \text{Erfolg} & 2/3\approx0{,}67 & 8\\
4 & 0+8(0{,}5)=4 & 16>0 & \text{Miss} & 2/4=0{,}5 & 16\\
5 & 0+16(-0{,}1)=-1{,}6 & 2{,}56>0 & \text{Miss} & 2/5=0{,}4 & 32\\
\end{array}
\]
Bester Punkt: $x=0$, $f=0$. (Nach Miss in It 1 wird $\sigma$ halbiert, danach mehrfach verdoppelt.)

\uebung{18.1 --- ES 1D}{mittel}
$f(x)=x^2$, Start $x=10$, $\sigma=2$, $z$-Folge $-2,\,-1,\,0{,}5,\,-0{,}1,\,0{,}1$. Führe 5 Iterationen aus
(Tabelle wie oben) und gib den besten Punkt sowie das finale $\sigma$ an.
\karo{9cm}

\uebung{18.2 --- ES 2D}{schwer}
$f(x_1,x_2)=x_1^2+x_2^2$, Start $(3,3)$, $\vec\sigma=(1,1)$. $z$-Folge (jeweils $(z_1,z_2)$):
$(-1,-2),(-0{,}5,0),(0{,}5,0{,}5),(-0{,}1,-0{,}1),(0{,}1,-0{,}1)$. Führe 5 Iterationen aus.
\karo{9.5cm}

% ============================================================
% THEMA 19 — Simulated Annealing
% ============================================================
\thema{19}{Metaheuristiken --- Simulated Annealing}
\bsp
Eine Verschlechterung um $\Delta f=2$ wird bei Temperatur $T=4$ vorgeschlagen; Zufallszahl $u=0{,}5$.
Akzeptiert?
\loesweg
Akzeptanzwahrscheinlichkeit (Metropolis): $p=\exp(-\Delta f/T)=\exp(-2/4)=\exp(-0{,}5)\approx0{,}607$.
Akzeptiere, falls $u<p$. Hier $u=0{,}5<0{,}607\Rightarrow$ \textbf{akzeptiert}.
\par\emph{Idee:} hohe $T\Rightarrow$ Verschlechterungen oft akzeptiert (globale Exploration); $T$ sinkt über den
Abkühlungsplan $\Rightarrow$ Verfahren wird gierig (lokale Feinsuche).

\uebung{19.1 --- Annahme prüfen}{leicht}
$\Delta f=3$, $T=2$, $u=0{,}4$. Berechne $p$ und entscheide, ob die schlechtere Lösung akzeptiert wird.
\karo{6cm}

\uebung{19.2 --- Einfluss der Temperatur}{mittel}
Eine Verschlechterung $\Delta f=1$ soll akzeptiert werden. Berechne $p$ für $T=4$ und für $T=0{,}5$ und
erkläre, was das über den Abkühlungsplan aussagt.
\karo{7cm}

% ============================================================
% THEMA 20 — PSO
% ============================================================
\thema{20}{Metaheuristiken --- Particle Swarm Optimization (PSO)}
\bsp
Geschwindigkeits-Update $v'=w\,v+c_1 r_1 (p_{\text{best}}-x)+c_2 r_2 (g_{\text{best}}-x)$, dann $x'=x+v'$.
Gegeben $x=2,\ v=1,\ p_{\text{best}}=1,\ g_{\text{best}}=0,\ w=0{,}5,\ c_1=c_2=1,\ r_1=r_2=1$.
\loesweg
\[
v'=0{,}5\cdot1+1\cdot1\cdot(1-2)+1\cdot1\cdot(0-2)=0{,}5-1-2=-2{,}5,\qquad x'=2+(-2{,}5)=-0{,}5 .
\]
Drei Anteile: \emph{Trägheit} $w v$, \emph{kognitiv} (Richtung eigener bester Pos.), \emph{sozial}
(Richtung global bester Pos.).

\uebung{20.1 --- Ein Update 1D}{mittel}
$x=4,\ v=-1,\ p_{\text{best}}=3,\ g_{\text{best}}=1,\ w=0{,}7,\ c_1=c_2=1{,}5,\ r_1=0{,}5,\ r_2=0{,}4$.
Berechne $v'$ und $x'$.
\karo{7cm}

\uebung{20.2 --- Ein Update 2D}{mittel}
$\vec x=(2,2),\ \vec v=(0,0),\ \vec p_{\text{best}}=(1,1),\ \vec g_{\text{best}}=(0,0),\ w=0{,}5,\ c_1=c_2=1,\
r_1=r_2=1$. Berechne $\vec v'$ und $\vec x'$ und benenne die drei Geschwindigkeitsanteile.
\karo{7.5cm}

% ============================================================
% THEMA 21 — ACO
% ============================================================
\thema{21}{Metaheuristiken --- Ant Colony Optimization (ACO)}
\bsp
Drei Kanten haben Pheromon $\tau=(1,3,2)$ und Sichtbarkeit $\eta=(4,2,1)$; Parameter $\alpha=\beta=1$.
Mit welcher Wahrscheinlichkeit wird jede Kante gewählt?
\loesweg
Gewicht $w_i=\tau_i^\alpha\cdot\eta_i^\beta$: $w_1=1\cdot4=4,\ w_2=3\cdot2=6,\ w_3=2\cdot1=2$, Summe $=12$.
\[
p_i=\frac{w_i}{\sum w}\ \Rightarrow\ p=\Big(\tfrac{4}{12},\tfrac{6}{12},\tfrac{2}{12}\Big)=\Big(\tfrac13,\tfrac12,\tfrac16\Big).
\]

\uebung{21.1 --- Pheromon-Update}{mittel}
Für drei Kanten gelte $\tau=(4,2,6)$, Verdunstungsrate $\rho=0{,}2$ und neu abgelegtes Pheromon
$\Delta\tau=(1,0,2)$. Berechne $\tau'=(1-\rho)\tau+\Delta\tau$.
\karo{6.5cm}

\uebung{21.2 --- Wahlwahrscheinlichkeit}{mittel}
Drei Kanten mit $\tau=(2,2,1)$, $\eta=(1,3,3)$, $\alpha=\beta=1$. Berechne die Auswahlwahrscheinlichkeiten.
\karo{6.5cm}

% ============================================================
% THEMA 22 — EA-Operatoren
% ============================================================
\thema{22}{Metaheuristiken --- EA-Operatoren (binär)}
\bsp
Elter $x=(1,0,1,1,0,1)$. (a) Bitflip-Mutation an Position 3. (b) Single-Point-Crossover von
$x_1=(1,1,0,0)$ und $x_2=(0,0,1,1)$ mit Schnitt nach Position 2.
\loesweg
\textbf{(a) Bitflip:} an Position 3 steht eine $1$; kippen $\Rightarrow 0$: $x^*=(1,0,\mathbf{0},1,0,1)$.
\par\textbf{(b) Crossover} (erste 2 Bits von $x_1$, Rest von $x_2$): $x_3=(1,1\,|\,1,1)=(1,1,1,1)$.

\uebung{22.1 --- Crossover + Mutation}{leicht}
$x_1=(1,0,1,1,0,1,0,1)$, $x_2=(0,1,0,0,1,0,1,0)$. (a) Single-Point-Crossover mit Schnitt nach Position 5
(Nachkomme $=$ vordere 5 Bits von $x_1$, Rest von $x_2$). (b) Bitflip von $x_1$ an Position 4.
\karo{6.5cm}

\uebung{22.2 --- Zweipunkt-Crossover}{mittel}
$x_1=(1,1,1,1,1,1)$, $x_2=(0,0,0,0,0,0)$. Führe einen \textbf{Zweipunkt}-Crossover mit Schnitten nach
Position 2 und nach Position 4 durch (Nachkomme übernimmt den Mittelteil von $x_2$, Außenteile von $x_1$).
\karo{6.5cm}

% ============================================================
% THEMA 23 — Dominanz & Pareto-Menge
% ============================================================
\thema{23}{Mehrziel --- Dominanz \& Pareto-Menge}
\bsp
Punkte im Zielraum (beide Ziele \textbf{minimieren}):
$P_1=(2,6),P_2=(3,4),P_3=(5,3),P_4=(4,5),P_5=(6,2)$. Bestimme die Pareto-Menge.
\loesweg
$\vec a$ \emph{dominiert} $\vec b$, wenn $\vec a$ in keinem Ziel schlechter und in mindestens einem besser ist.
\par$P_2=(3,4)$ dominiert $P_4=(4,5)$ ($3<4$ und $4<5$). Sonst dominiert kein Punkt einen anderen.
\par\emph{Pareto-optimal} (von keinem dominiert): $\boxed{\{P_1,P_2,P_3,P_5\}}$; $P_4$ ist dominiert.

\uebung{23.1 --- Pareto-Menge 2D}{mittel}
$P_1=(1,5),P_2=(2,3),P_3=(4,2),P_4=(3,4),P_5=(5,1),P_6=(2,6)$ (beide minimieren). Bestimme alle
dominierten Punkte und die Pareto-Menge.
\karo{7.5cm}

\uebung{23.2 --- Drei Ziele}{mittel}
$P_1=(1,2,3),P_2=(2,1,2),P_3=(3,2,1),P_4=(2,2,2),P_5=(1,1,1)$ (alle minimieren). Bestimme die Pareto-Menge.
\karo{7cm}

% ============================================================
% THEMA 24 — Non-Dominated Sorting
% ============================================================
\thema{24}{Mehrziel --- Non-Dominated-Sorting (Ränge)}
\bsp
$P_1=(1,2,3),P_2=(2,1,2),P_3=(3,2,1),P_4=(2,2,2),P_5=(3,3,3),P_6=(1,3,1)$ (alle minimieren). Bestimme die
Ränge.
\loesweg
\textbf{Rang 1} = nicht dominierte Punkte: $P_4$ wird von $P_2$ dominiert ($2,1,2\le2,2,2$), $P_5$ von vielen.
$P_1,P_2,P_3,P_6$ werden von keinem dominiert $\Rightarrow$ Rang 1 $=\{P_1,P_2,P_3,P_6\}$.
\par Diese entfernen; unter $\{P_4,P_5\}$ ist $P_4$ nicht dominiert $\Rightarrow$ \textbf{Rang 2} $=\{P_4\}$.
\par Übrig $P_5\Rightarrow$ \textbf{Rang 3} $=\{P_5\}$.

\uebung{24.1 --- Ränge 2D (I)}{mittel}
$P_1=(1,4),P_2=(2,2),P_3=(3,1),P_4=(2,5),P_5=(4,3),P_6=(5,2)$ (minimieren). Bestimme alle Ränge.
\karo{8cm}

\uebung{24.2 --- Ränge 2D (II)}{mittel}
$P_1=(1,1),P_2=(2,3),P_3=(3,2),P_4=(4,4),P_5=(2,2)$ (minimieren). Bestimme alle Ränge.
\karo{8cm}

% ============================================================
% THEMA 25 — Crowding Distance + NSGA-II
% ============================================================
\thema{25}{Mehrziel --- Crowding Distance \& NSGA-II-Auswahl}
\bsp
Nichtdominierte Front (minimieren): $A=(1,7),B=(2,5),C=(3,4),D=(5,3),E=(7,1)$. Berechne die Crowding
Distance und wähle 4 Punkte aus.
\loesweg
Randpunkte erhalten $\infty$: $A,E$. Spannweiten $f_1{:}\,7-1=6$, $f_2{:}\,7-1=6$. Für innere Punkte
$CD=\sum_{\text{Ziel}}\frac{\text{Nachbar}_{\text{rechts}}-\text{Nachbar}_{\text{links}}}{\text{Spannweite}}$
(nach $f_1$ sortiert):
\[
CD(B)=\tfrac{3-1}{6}+\tfrac{7-4}{6}=\tfrac{5}{6}\approx0{,}83,\quad
CD(C)=\tfrac{5-2}{6}+\tfrac{5-3}{6}=\tfrac{5}{6}\approx0{,}83,\quad
CD(D)=\tfrac{7-3}{6}+\tfrac{4-1}{6}=\tfrac{7}{6}\approx1{,}17 .
\]
\textbf{Auswahl} (4 Punkte): zuerst Ränder $A,E$ ($\infty$), dann größte $CD$: $D$, dann $B$ (oder $C$).
$\Rightarrow\{A,E,D,B\}$.

\uebung{25.1 --- Crowding Distance}{schwer}
Front (minimieren): $F=(1,9),G=(3,6),H=(4,4),I=(6,3),J=(9,1)$. Berechne die Crowding Distance aller Punkte
und wähle 4 aus.
\karo{9cm}

\uebung{25.2 --- NSGA-II-Auswahl}{mittel}
Eine Population: $a$ (Rang 1, $CD=\infty$), $b$ (Rang 1, $CD=0{,}5$), $c$ (Rang 1, $CD=1{,}2$),
$d$ (Rang 2, $CD=\infty$), $e$ (Rang 1, $CD=0{,}8$), $f$ (Rang 2, $CD=1{,}0$). Es sollen \textbf{3}
Individuen überleben. Welche werden ausgewählt? Begründe.
\karo{7cm}

% ============================================================
% THEMA 26 — Hypervolumen + HVC
% ============================================================
\thema{26}{Mehrziel --- Hypervolumen \& HV-Beitrag}
\bsp
Front (minimieren): $P_1=(2,8),P_2=(4,5),P_3=(7,3)$, Referenzpunkt $\vec r=(10,10)$. Berechne das
Hypervolumen und die HV-Beiträge.
\loesweg
Nach $f_1$ sortieren, Treppenflächen $\sum (f_{1,\text{next}}-f_{1,i})\,(r_2-f_{2,i})$:
\[
P_1{:}(4-2)(10-8)=4,\quad P_2{:}(7-4)(10-5)=15,\quad P_3{:}(10-7)(10-3)=21\ \Rightarrow\ HV=40 .
\]
$HVC(P_i)=HV-HV(\text{ohne }P_i)$:
\[
HVC(P_1)=40-36=4,\quad HVC(P_2)=40-31=9,\quad HVC(P_3)=40-34=6 .
\]

\uebung{26.1 --- HV \& HVC (I)}{schwer}
Front (minimieren): $(1,6),(3,4),(5,2)$, Referenz $\vec r=(8,8)$. Berechne $HV$ und alle HV-Beiträge.
\karo{9cm}

\uebung{26.2 --- HV \& HVC (II)}{schwer}
Front (minimieren): $(2,7),(3,5),(6,2)$, Referenz $\vec r=(8,8)$. Berechne $HV$ und alle HV-Beiträge.
\karo{9cm}

% ============================================================
% THEMA 27 — Simplex-Tableau
% ============================================================
\thema{27}{Lineare Optimierung --- Simplex-Tableau}
\bsp
$\max z=3x_1+2x_2$ unter $x_1+x_2\le 4$, $x_1+3x_2\le 6$, $x_1,x_2\ge0$ (vgl.\ Thema 5, jetzt rechnerisch).
\loesweg
\textbf{(1) Standardform} mit Slacks: $x_1+x_2+s_1=4$, $x_1+3x_2+s_2=6$. Start-Basis $(s_1,s_2)$,
$z$-Zeile $z-3x_1-2x_2=0$.
\[
\begin{array}{c|cccc|c}
 & x_1 & x_2 & s_1 & s_2 & b\\\hline
s_1 & 1 & 1 & 1 & 0 & 4\\
s_2 & 1 & 3 & 0 & 1 & 6\\\hline
z & -3 & -2 & 0 & 0 & 0
\end{array}
\]
\textbf{(2) Pivot:} negativster $z$-Eintrag $-3\Rightarrow x_1$ kommt in die Basis. Quotiententest
$\tfrac{4}{1}=4,\ \tfrac{6}{1}=6\Rightarrow$ kleinster $=4\Rightarrow s_1$ verlässt; Pivotelement $1$.
Zeilenumformungen (Pivotzeile bleibt, andere eliminieren):
\[
\begin{array}{c|cccc|c}
 & x_1 & x_2 & s_1 & s_2 & b\\\hline
x_1 & 1 & 1 & 1 & 0 & 4\\
s_2 & 0 & 2 & -1 & 1 & 2\\\hline
z & 0 & 1 & 3 & 0 & 12
\end{array}
\]
\textbf{(3)} $z$-Zeile $\ge0\Rightarrow$ \textbf{optimal}: $x_1=4,\ x_2=0,\ z=12$ (identisch zum grafischen Ergebnis).

\uebung{27.1 --- Ein Pivot}{mittel}
$\max z=3x_1+x_2$ unter $x_1+x_2\le 4$, $2x_1+x_2\le 6$, $x_1,x_2\ge0$. Löse mit dem Simplex-Tableau
(Standardform, Pivot mit Quotiententest) und gib das Optimum an.
\karo{9.5cm}

\uebung{27.2 --- Zwei Pivots}{schwer}
$\max z=2x_1+3x_2$ unter $x_1+x_2\le 4$, $x_1+2x_2\le 6$, $x_1,x_2\ge0$. Löse mit dem Simplex-Tableau
(zwei Pivotschritte nötig).
\karo{10cm}

% ============================================================
% THEMA 28 — Quadratische Optimierung
% ============================================================
\thema{28}{Quadratische Optimierung (QO)}
\bsp
Ist $\min f(x,y)=x^2+2y^2+xy-4x$ unter $x+y\le 2$, $x,y\ge0$ ein QO-Problem? Ist es konvex?
\loesweg
Ein \textbf{QO-Problem} hat eine \emph{quadratische} Zielfunktion und \emph{lineare} Nebenbedingungen,
also Form $\min\tfrac12\vec x^\top \mathbf Q\vec x+\vec c^\top\vec x$ u.\,d.\,N.\ $\mathbf A\vec x\le \vec b$.
\par Hier: $f$ ist quadratisch mit $\mathbf Q=\begin{pmatrix}2&1\\1&4\end{pmatrix}$ (denn
$\tfrac12\vec x^\top\mathbf Q\vec x=x^2+2y^2+xy$) und $\vec c=(-4,0)$; alle Nebenbedingungen sind linear
$\Rightarrow$ \textbf{ja, QO-Problem}. $\mathbf Q$: Hauptminoren $2>0,\ \det=7>0\Rightarrow$ positiv definit
$\Rightarrow$ \textbf{konvexes} QO-Problem (jedes lokale Optimum ist global).

\uebung{28.1 --- QO oder nicht?}{mittel}
Ist $\min f(x,y)=x^2+y^2$ unter $x^2+y^2\le 1$, $x+y\ge 1$ ein QO-Problem? Begründe (welche Bedingung
verletzt die QO-Definition?).
\karo{6.5cm}

\uebung{28.2 --- In Matrixform bringen}{mittel}
Schreibe $f(x,y)=2x^2+3y^2-xy+x-2y$ in der Form $\tfrac12\vec x^\top\mathbf Q\vec x+\vec c^\top\vec x$
(also $\mathbf Q$ und $\vec c$ angeben) und prüfe $\mathbf Q$ auf Definitheit (Konvexität).
\karo{7.5cm}

% ============================================================
% THEMA 29 — SMBO / Surrogat-Modelle
% ============================================================
\thema{29}{Surrogat-Modell-basierte Optimierung (SMBO)}
\bsp
Wozu dient SMBO und wie läuft eine Iteration ab?
\loesweg
SMBO wird für \textbf{teure} Zielfunktionen (lange Simulation, Experiment) eingesetzt: Statt das echte $f$
oft auszuwerten, baut man ein \emph{billiges Surrogat-Modell} (z.\,B.\ Gauß-Prozess/Kriging, Polynom) aus
den bisherigen Auswertungen. \textbf{Iteration:}
(1) Surrogat an die Datenpunkte anpassen;
(2) den als Nächstes auszuwertenden Punkt über das Modell bestimmen (z.\,B.\ \emph{Minimum des Surrogats}
bzw.\ eine Acquisition-Funktion, die auch Unsicherheit berücksichtigt);
(3) dort das \emph{echte} (teure) $f$ auswerten und den neuen Punkt aufnehmen;
(4) wiederholen, bis das Budget erschöpft ist. So nähert man sich mit wenigen echten Auswertungen dem Optimum.

\uebung{29.1 --- Nächster Auswertungspunkt}{mittel}
Ein Surrogat-Modell hat sein Minimum bei $x\approx 1{,}2$; das wahre (unbekannte) Minimum liegt bei
$x\approx 0{,}5$. (a) Wo wertet SMBO als Nächstes das \emph{echte} $f$ aus? (b) Was passiert mit dem Modell
und seinem Minimum nach dieser Auswertung?
\karo{6.5cm}

\uebung{29.2 --- Wann SMBO?}{leicht}
Nenne zwei Eigenschaften eines Problems, bei denen SMBO besonders sinnvoll ist, und erkläre, welche Rolle
das Surrogat-Modell dabei spielt.
\karo{6cm}

% ============================================================
% THEMA 30 — Verfahrensüberblick & CMA-ES
% ============================================================
\thema{30}{Verfahrensüberblick \& CMA-ES (Einordnung)}
\bsp
Ordne die Verfahren nach den Eigenschaften \emph{global?}, \emph{gradientenbasiert?},
\emph{deterministisch?}, \emph{\#Ziele} ein und erkläre, wodurch sich CMA-ES von der $(1+1)$-ES unterscheidet.
\loesweg
\[
\renewcommand{\arraystretch}{1.2}
\begin{array}{l|c|c|c|c}
\text{Verfahren} & \text{global} & \text{gradientenbasiert} & \text{deterministisch} & \text{\#Ziele}\\\hline
\text{Stat.\ Punkte/Hesse} & \checkmark & \checkmark & \checkmark & 1\\
\text{Gradientenabstieg} & \times & \checkmark & \checkmark & 1\\
\text{Simplex (LP)} & \times & \checkmark & \checkmark & 1\\
\text{KKT} & \checkmark & \checkmark & \checkmark & 1\\
\text{Nelder-Mead} & \times & \times & \checkmark & 1\\
\text{DIRECT} & \checkmark & \times & \checkmark & 1\\
(1{+}1)\text{-ES} & \times & \times & \times & 1\\
\text{NSGA-II} & \checkmark & \times & \times & >1
\end{array}
\]
\textbf{CMA-ES vs.\ $(1+1)$-ES:} Beide sind stochastische Evolutionsstrategien. CMA-ES arbeitet mit einer
\emph{Population} und passt die \emph{komplette Kovarianzmatrix} der Mutationsverteilung selbstadaptiv an
(Schrittweite \emph{und} Richtungen), wodurch es deutlich \emph{globaler} und robuster ist; die $(1+1)$-ES
hat nur ein Elter/Kind und steuert allein die Schrittweite $\sigma$ über die 1/5-Erfolgsregel.

\uebung{30.1 --- Verfahren zuordnen}{mittel}
Ordne ein (global? / gradientenbasiert? / deterministisch? / \#Ziele): \emph{Simulated Annealing},
\emph{PSO}, \emph{Newton-Verfahren}, \emph{SMS-EMOA}.
\karo{7.5cm}

\uebung{30.2 --- Verfahren auswählen}{mittel}
Begründe kurz, welches Verfahren du jeweils wählst:
\begin{enumerate}[label=(\alph*),nosep]
  \item eine \emph{teure} Black-Box-Funktion mit nur wenigen erlaubten Auswertungen;
  \item ein \emph{lineares} Problem mit linearen Nebenbedingungen;
  \item ein Problem mit \emph{zwei} konkurrierenden Zielen.
\end{enumerate}
\karo{7cm}

\end{document}
